\section{Síntesis y ampliación}

En esta sesión se partió de las dos variantes de Qwen3, Dense y MoE, y mediante la comparación de configuraciones y el cálculo del número de parámetros paso por paso se construyó una comprensión intuitiva de la arquitectura Sparse MoE:

\begin{enumerate}
  \item MoE sustituye la FFN por múltiples Expert independientes más un Router, lo que logra una activación dispersa
  \item GQA comprime la KV Cache compartiendo cabezas KV, y es hoy el estándar de facto en los LLM
  \item Qwen3 rompe la restricción tradicional $d_{\text{model}} = n_{\text{heads}} \times d_{\text{head}}$
  \item De los 30B de parámetros totales, en inferencia solo se activan alrededor de 3.3B, con una eficiencia de inferencia cercana a la de un modelo Dense del mismo tamaño
\end{enumerate}

\subsection{Lecturas adicionales}

\begin{itemize}
  \item El blog de Sebastien Raschka: serie de explicaciones detalladas sobre arquitecturas de LLM
  \item El informe técnico oficial de Qwen3
  \item El diseño de MoE en los informes técnicos de DeepSeek V2/V3
  \item Las explicaciones sobre LLaMA 2 y GQA en la serie «Post NOS Chat GPT»
\end{itemize}

%% --- Fin del contenido --- %%

\end{document}
